\section{Projective content and a boundary to the single-map inverse}
\label{pc-section}

The rational maps of Section~\ref{rd-section} send every actual profile
to a rational point. A converse must account for integral content.
That content admits a bound independent of the exponent, but the bound
alone does not preserve the prescribed exponent after primitive reduction.
The following counterfamily concerns a single power map; it does not
assert a rational point on the full double fiber product.

Retain $\mathcal O=\mathbb Z[\zeta]$, $\zeta^2-\zeta+1=0$.
For $z=A+B\zeta\ne0$, write
$\operatorname{cont}(z)=\gcd(|A|,|B|)$.

\begin{lemma}[A content divisor]\label{pc-content}
If $0\ne\tau\in\mathcal O$ and $W\in\mathcal O$ is primitive, then
\[
 \operatorname{cont}(\tau W)\mid\operatorname N\tau.
\]
\end{lemma}
\begin{proof}
Write $\tau=r+s\zeta$, $W=u+v\zeta$ and $\tau W=A+B\zeta$.
Multiplication and its adjugate give
\[
 A=ru-sv,\quad B=su+(r+s)v,
\]
\[
 (r+s)A+sB=(\operatorname N\tau)u,\qquad
 -sA+rB=(\operatorname N\tau)v.
\]
The content divides both right sides. Bezout for $\gcd(|u|,|v|)=1$
proves the claim. No coprimality between $\tau$ and $W$ is needed.
\end{proof}

For an unramified primitive root $w$, its powers are primitive by the
reviewed power theorem. Thus the content $C$ of $\tau w^n$ divides
$\operatorname N\tau$, independently of $n$. If $(a,b)$ is the
primitive output, however, its norm is
\begin{equation}\label{pc-reduced-norm}
 a^2+ab+b^2=\frac{\operatorname N\tau}{C^2}
                         (\operatorname Nw)^n.
\end{equation}
The residual displayed here is rational and need not be integral.

\begin{theorem}[Bounded content does not preserve the specified exponent]
\label{pc-inverse-boundary}
Put $\pi=2+\zeta$ and $w=\bar\pi=3-\zeta$.
For every integer $n\ge2$ there is a unit $u_n$ such that, with
$\tau_n=u_n\pi$, the primitive output of $\tau_nw^n$ has positive
coordinates $a_n,b_n$ and
\begin{gather}
 \gcd(a_n,b_n)=1,\qquad a_n^2+a_nb_n+b_n^2=7^{n-1},
 \label{pc-positive-norm}\\
 \operatorname{cont}(\tau_nw^n)=7,
 \qquad\operatorname N\tau_n=\operatorname Nw=7.\notag
\end{gather}
There are no positive integers $V,R$ with $R>1$ satisfying
\[
 a_n^2+a_nb_n+b_n^2=VR^n.
\]
The multipliers $\tau_n$ belong to a fixed set of six elements, and
the input parameter in the rational power map is always $t=-3$.
\end{theorem}
\begin{proof}
The identity $\pi\bar\pi=7$ gives
\[
 \pi w^n=7w^{n-1}.
\]
Every power $w^m$, $m\ge1$, is primitive. Indeed, its norm is $7^m$,
so any common rational prime divisor of its coordinates would be seven.
Evaluation at $\zeta=5$ defines a homomorphism
$\mathcal O\to\mathbb F_7$, since $5^2-5+1=0$ modulo seven.
It sends $w$ to $-2\ne0$, so $w^m$ cannot have both coordinates
divisible by seven.

No boundary arm of $w^m$ vanishes: for a primitive element this would
force norm one. A unit sends it into the open positive sector. To see
all sign cases explicitly, multiplication by $\zeta$ and $\zeta^2$
sends $(A,B)$ to $(-B,A+B)$ and $(-A-B,A)$. If the signs are opposite,
negate first if needed to have $A>0>B$; the two transformations work
respectively when $A+B>0$ and when $A+B<0$. Equal signs require only
the identity or negation. Choose $u_n$ for $w^{n-1}$ in this way.
Units preserve content and norm, proving \eqref{pc-positive-norm}.

If $VR^n=7^{n-1}$ with $R>1$, every prime divisor of $R$ is seven.
Thus $R=7^j$ for some $j\ge1$, and the valuation of $VR^n$ at seven
is at least $n$, a contradiction. Finally the root $w=3-\zeta$
has rational parameter $3/(-1)=-3$. Replacing it by $t+\zeta=-w$
only scales the output by $(-1)^n$, so the projective output is unchanged.
\end{proof}

An explicit positive instance is $n=5$, $u_5=\zeta^2$. Then
\[
 \tau_5=-3+2\zeta,\qquad w^4=39-55\zeta,
\]
\[
 \tau_5w^5=112+273\zeta=7(16+39\zeta),\qquad
 16^2+16\cdot39+39^2=2401=7^4.
\]

In this family the canonical $n$-free residual of the reduced norm
is $7^{n-1}$ and its canonical $n$-th-power root is one. Thus its
normalized logarithmic residual tends to $\log7$, although the
given multiplier norm has normalized logarithm $\log7/n\to0$.
The bounded cancellation leaves the rational residual $1/7$ in
\eqref{pc-reduced-norm}; integral normalization then changes the
residual exponent by almost $n$.

This excludes only an automatic single-map inverse with the same
prescribed exponent and an integral nonunit root. The evident alternative
exponent $n-1$ remains available. The family is not asserted to meet
the second map of the RD fiber product, and does not refute the forward
DC lift, simultaneous compatibility, signed compensation, or $abc$.
The statements here are ordinary proofs; no Lean verification is claimed.
